\section{Experimental Evaluation}
\label{sec:experimental_evaluation}
\label{sec:experiments}

\subsection{Data and Evaluation Protocol}
\label{sec:data_protocol}

The evaluation uses a real-world MBB operational dataset provided by Huawei Technologies. The source collection covers approximately 500 anonymized cellular sites, while the aligned panel evaluated in this study contains 33 anonymized cell entities after task-specific panel construction. The processed panel comprises 103,983 cell-time instances, 3,151 unique hourly timestamps, 48 KPI/business counters and 137 alarm/fault columns. Only 11 alarm/fault columns are active, and 1.62\% of rows contain an active alarm. The panel contains one 21 days 01:00:00 time discontinuity; forecasting windows crossing this gap are removed. Dataset dimensions and feature sparsity characteristics are detailed in Table~\ref{tab:data_stats}.

All experiments use chronological splitting with 70\% training, 15\% validation and 15\% testing. Feature scalers, degradation thresholds, inferred cell interaction graphs and event influence priors are fitted on the training partition only. Models are evaluated at four horizons ($h \in \{1, 3, 6, 12\}$ hours). The primary task is continuous future KPI value forecasting. Secondary tasks are weak KPI degradation classification based on train-derived threshold violations and weak alarm-candidate prioritization based on event intensity and lift.

\begin{table}[t]
\centering
\caption{Dataset statistics for the anonymized Huawei MBB operations data and evaluated cell-time panel.}
\label{tab:data_stats}
\begin{tabularx}{\textwidth}{lX}
\toprule
Item & Value \\
\midrule
Collection scope & Approximately 500 MBB cellular sites in the source collection \\
Rows & 103,983 \\
Original columns & 187 \\
Anonymized cell entities in evaluated panel & 33 \\
Unique hourly timestamps & 3,151 \\
Observed time structure & 3,151 unique hourly timestamps in the evaluated panel; exact site-identifiable information is not used \\
KPI/business columns & 48 \\
Alarm/fault columns & 137 \\
Active alarm/fault columns & 11 \\
Rows with active alarms & 1,681 (1.62\%) \\
Missing values & None detected \\
Detected time discontinuity & One 21 days 01:00:00 gap; crossing windows removed \\
\bottomrule
\end{tabularx}

\end{table}


\subsection{Compared Models and Metrics}
\label{sec:models_metrics}

The runnable comparison includes Ridge regression for forecasting, a linear logistic classifier for weak degradation classification, corresponding no-alarm variants, recency-weighted event-intensity baselines, event-feature scoring baselines, CET-STGL-sklearn, its ablations, and two fixed neural sequence baselines: Long Short-Term Memory (LSTM) and Temporal Convolutional Network (TCN). "Ridge/Logistic" is used as a compact cross-task label in the result tables: forecasting uses Ridge regression, whereas classification uses the SGD-trained linear logistic classifier. LSTM and TCN use the same 24-hour observation window, chronological split, train-only scaling, active alarm inputs and fixed hyperparameters without search. They are sequence-model comparators rather than full neural spatio-temporal graph implementations.

Continuous KPI forecasting is measured using Mean Absolute Error (MAE), Root Mean Square Error (RMSE), and Mean Absolute Percentage Error (MAPE) with a numerical floor $\max(|y|, 1.0)$ to stabilize zero-value division. Degradation classification performance is quantified via Area Under the Precision-Recall Curve (AUPRC), Area Under the ROC Curve (AUROC), F1-score, Precision, and Recall. Weak alarm-candidate ranking consistency is evaluated using the hit rate at cutoff $K$, denoted $\mathrm{Hit}(K)$ for $K \in \{1, 3, 5\}$, Mean Reciprocal Rank (MRR), and Normalized Discounted Cumulative Gain at cutoff 5, denoted $\mathrm{nDCG}(5)$. Here, $\mathrm{Hit}(K)$ means that the weak target alarm appears among the top-$K$ ranked candidates, and $\mathrm{nDCG}(5)$ is computed after truncating the ranked list at five candidates.
